\documentclass{article}
\usepackage{graphicx}
\usepackage{polski}
\usepackage{amsmath}
\usepackage{indentfirst}
\usepackage{natbib}
\usepackage{hyperref}
\usepackage[utf8]{inputenc}
\usepackage{listings}
\lstset{language=C++,
                basicstyle=\ttfamily,
}
\usepackage{anysize}
\marginsize{3cm}{2.5cm}{2.5cm}{2.5cm}%LPGD
\begin{document}
\begin{center}
POLITECHNIKA ŚLĄSKA\\
WYDZIAŁ INŻYNIERII BIOMEDYCZNEJ\\
\vspace{\stretch{2}}
\begin{figure}[ht]
\centering
\includegraphics[width=0.4\textwidth]{rys/logo.png}
\end{figure}
\textbf{Nazwa przedmiotu}\\
\vspace{\stretch{1}}
\huge{\textbf{Nazwa projektu}}\\
\vspace{\stretch{1}}
\large{Autor}\\
\vspace{\stretch{5}}
\large{\today}
\end{center}
\thispagestyle{empty}
\newpage

\tableofcontents
\newpage

\listoffigures

\newpage

\section{Wstęp}

\subsection{Cel projektu}

\section{Metodologia}

\subsection{Wczytywanie obrazów}

\begin{figure}
\centering
\includegraphics[width=1\textwidth]{rys/schematblokowy.png}
\caption{Schemat poszczególnych etapów algorytmu wraz z użytymi metodami biblioteki ITK}
\label{fig:schemat}
\end{figure}

Do obliczania współczynnika DICE, używano masek z bazy \citep{mri}, a w niektórych przypadkach wykonano je ręcznie przy pomocy oprogramowania SLICER używając metody \textit{Grow from seeds}. Przykładowa ręcznie wykonana maska (przekrój) został przedstawiony na rysynku (Rys. \ref{fig:schemat}).

\section{Specyfikacja wewnętrzna}

\section{Specyfikacja zewnętrzna}

\section{Wyniki}

\section{Podsumowanie}

\bibliographystyle{plain}
\bibliography{references}

\end{document}